% =============================================================================
% Worked solutions for Listas/PS2.pdf (Corporate Finance -- Problem Set 2).
% Short answer key: Listas/PS2-solutions-short.tex.
% Compile: python -I Listas/PS2-solutions_check.py (writes PS2-solutions-table.tex), then
%          pdflatex (twice); check: pdffonts (Type 1 only, uni = yes).
% =============================================================================
\documentclass[11pt,a4paper]{article}
\usepackage[utf8]{inputenc}
\usepackage[T1]{fontenc}
\usepackage{lmodern}   % Type 1 vector fonts (replaces PK bitmaps)
\usepackage{cmap}      % ToUnicode CMap -> copyable text and accents
\usepackage[english]{babel}
\usepackage{amsmath,amssymb,amsfonts,amsthm}
\usepackage[a4paper,margin=2.2cm]{geometry}
\usepackage{booktabs}
\usepackage{longtable}
\usepackage{enumitem}
\usepackage{fancyhdr}
\usepackage{xcolor}
\usepackage{microtype}

\pagestyle{fancy}\fancyhf{}
\lhead{\footnotesize Finan\c{c}as II -- Corporativas (MPE Insper)}
\rhead{\footnotesize Problem Set 2 --- Worked solutions}
\cfoot{\footnotesize\thepage}
\renewcommand{\headrulewidth}{0.3pt}
\setlength{\headheight}{14pt}

\makeatletter
\renewcommand{\@seccntformat}[1]{\csname the#1\endcsname.\quad}
\makeatother
\numberwithin{equation}{section}
\setlist{topsep=2pt,itemsep=1pt}

\newcommand{\step}[2]{\par\medskip\noindent\textbf{Step #1 --- #2.}~}
\newenvironment{notation}{\par\noindent\textbf{Notation.}\begin{itemize}[leftmargin=1.4em]}{\end{itemize}}
\newcommand{\intuition}[1]{\par\medskip\noindent\textbf{Intuition.}~#1}
\newcommand{\refline}[1]{\par\smallskip\noindent{\small\color{black!70}\textbf{Ref:} #1}}
\newcommand{\remark}[1]{\par\smallskip\noindent\textit{Remark.}~#1}
\newcommand{\answer}[1]{\par\medskip\noindent\fbox{\parbox{\dimexpr\linewidth-2\fboxsep-2\fboxrule}{\textbf{Answer.}~#1}}\par}

\begin{document}

\begin{center}
  {\Large\bfseries Corporate Finance --- Problem Set 2}\\[4pt]
  {\large Worked solutions}
\end{center}

\noindent Statement: \texttt{Listas/PS2.pdf}. Each exercise is restated briefly, then solved
step by step. Rates are written as decimals inside equations. References are to the lecture
slides (\texttt{Aula/Topic2A.pdf}, \texttt{Aula/Topic2B.pdf}) and to Berk, DeMarzo \& Harford,
\emph{Fundamentals of Corporate Finance}, 2nd ed.\ (2011). Every number below is reproduced by
\texttt{Listas/PS2-solutions\_check.py}, which also generates the table of Appendix~A. A short
answer key is in \texttt{Listas/PS2-solutions-short.pdf}. The statement labels two parts of
Exercise~1 as ``(e)''; the second is called (e$'$) here.

% =============================================================================
\section{Exercise 1 --- State prices, options and MM in a complete market}

\noindent\textbf{Data.} Two equally likely states at $t=1$, slump ($s$) and boom ($b$). Stock: price
$1.05$, payoff $(0,\ 2.5)$. Bond: price $0.92$, payoff $(1,\ 1)$. Project: investment $1$ at $t=0$,
payoff $(0.5,\ 3)$. MM assumptions hold, no taxes. The entrepreneur has no money.

\begin{notation}
  \item $\psi_s,\psi_b>0$: state prices, the $t=0$ price of 1 paid only in the slump (boom).
  \item $r_f$: risk-free rate; $q_s,q_b$: risk-neutral probabilities, $q_\omega=\psi_\omega(1+r_f)$.
  \item $X=(X_s,X_b)=(0.5,\,3)$: project payoff; $A$: market value of the project (the firm's assets).
  \item $C(K)$, $P(K)$: prices of a call and a put on the project with strike $K\ge0$.
  \item $D$: amount raised with debt (= market value of debt); $F$: face value of debt.
  \item $E$: market value of equity; $r_A,r_D,r_E$: expected returns, computed with the true
        probabilities $\tfrac12,\tfrac12$.
\end{notation}

\subsection*{(a) State prices}

\step{1}{Pricing equations} A traded asset's price is the sum of its payoffs times state prices:
\begin{align}
  \text{Stock:}\quad 0\,\psi_s + 2.5\,\psi_b &= 1.05, \label{eq:1stock}\\
  \text{Bond:}\quad 1\,\psi_s + 1\,\psi_b &= 0.92. \label{eq:1bond}
\end{align}

\step{2}{Solve} From \eqref{eq:1stock}:
\begin{equation}\label{eq:1psib}
  \psi_b = \frac{1.05}{2.5} = 0.42 .
\end{equation}
Substituting \eqref{eq:1psib} into \eqref{eq:1bond}:
\begin{equation}\label{eq:1psis}
  \psi_s + (0.42) = 0.92 \quad\Longrightarrow\quad \psi_s = 0.50 .
\end{equation}

\step{3}{Implied rates and probabilities} The bond pays 1 for sure, so
$1+r_f = 1/0.92$, i.e.\ $r_f = 0.086957$ (8.70\%). The risk-neutral probabilities are
$q_s = 0.50/0.92 = 0.543478$ and $q_b = 0.42/0.92 = 0.456522$.

\answer{The slump claim costs $\psi_s=0.50$ and the boom claim $\psi_b=0.42$.}

\intuition{The two states are equally likely, yet a unit in the slump costs more than a unit in the
boom ($0.50>0.42$). Investors value consumption most when they are poor; insurance against the
slump is expensive. Equivalently, the risk-neutral probability of the slump (54.3\%) exceeds its
true probability (50\%).}

\step{4}{Value of the project} Using \eqref{eq:1psib}--\eqref{eq:1psis}:
\begin{align}
  A &= (0.5)(0.50) + (3)(0.42) \notag\\
    &= 0.25 + 1.26 = 1.51 . \label{eq:1A}
\end{align}
The project costs 1 and is worth 1.51: $\mathrm{NPV}=0.51$.

\subsection*{(b) Calls on the project}

A call with strike $K$ pays $\max(X_\omega-K,0)$ in state $\omega$.
\begin{align}
  C(1)   &= \max(0.5-1,0)(0.50) + \max(3-1,0)(0.42) = 0 + (2)(0.42) = 0.84, \label{eq:1C1}\\
  C(1.5) &= \max(0.5-1.5,0)(0.50) + \max(3-1.5,0)(0.42) = 0 + (1.5)(0.42) = 0.63. \label{eq:1C15}
\end{align}

\subsection*{(c) Puts on the project}

A put with strike $K$ pays $\max(K-X_\omega,0)$.
\begin{align}
  P(1)   &= \max(1-0.5,0)(0.50) + \max(1-3,0)(0.42) = (0.5)(0.50) + 0 = 0.25, \label{eq:1P1}\\
  P(1.5) &= \max(1.5-0.5,0)(0.50) + \max(1.5-3,0)(0.42) = (1)(0.50) + 0 = 0.50. \label{eq:1P15}
\end{align}

\subsection*{(d) Put--call parity and MM}

\step{1}{The middle term} $K/(1/0.92) = 0.92K$ is the price of a safe bond paying $K$ at $t=1$,
since $1/0.92 = 1+r_f$.

\step{2}{Substitute \eqref{eq:1C1}--\eqref{eq:1P15}}
\begin{align}
  C(1) + 0.92(1) - P(1)       &= (0.84) + 0.92 - (0.25) = 1.51, \label{eq:1lhs}\\
  C(1.5) + 0.92(1.5) - P(1.5) &= (0.63) + 1.38 - (0.50) = 1.51. \label{eq:1rhs}
\end{align}
Both sides equal $A$ from \eqref{eq:1A}.

\step{3}{Why it holds for every $K$} State by state, for any payoff $X_\omega$ and strike $K$:
\begin{equation}\label{eq:1parity}
  \max(X_\omega-K,0) + K - \max(K-X_\omega,0) = X_\omega .
\end{equation}
(If $X_\omega\ge K$: $(X_\omega-K)+K-0$; if $X_\omega<K$: $0+K-(K-X_\omega)$.) Pricing both sides with
state prices gives $C(K) + K/(1+r_f) - P(K) = A$ for every $K$: put--call parity.

\answer{Both sides equal $1.51=A$, the value of the project, for any strike.}

\intuition{Read $K$ as the face value of debt. With limited liability, equity pays
$\max(X-K,0)$: it is a call on the assets, worth $C(K)$. Risky debt pays $\min(X,K) = K - \max(K-X,0)$:
a safe bond minus a put, worth $K/(1+r_f)-P(K)$. The put is the creditors' loss from limited liability.
By \eqref{eq:1parity} equity plus debt is worth $A$ whatever $K$ is. A higher face value moves value
from the call to the bond, but the total does not change. This is MM~I, proved by put--call parity.}

\refline{Slides Topic 2A (complete markets, state prices, MM I with risky debt); Berk
\S21.5--21.6 (put--call parity; equity as a call), \S16.1--16.2.}

\subsection*{(e) All debt: raise 1 with debt}

\step{1}{Is the debt safe?} If $F\le0.5$ the debt pays $F$ in both states and is worth $0.92F$.
Raising 1 would need $F = 1/0.92 = 1.087 > 0.5$, a contradiction. So the debt defaults in the slump:
it pays $\min(X,F) = (0.5,\ F)$ with $0.5<F\le3$.

\step{2}{Face value} Debt must be worth what it raises:
\begin{equation}
  (0.5)(0.50) + F(0.42) = 1 \quad\Longrightarrow\quad F = \frac{1-0.25}{0.42} = 1.785714. \label{eq:1F}
\end{equation}

\step{3}{Values} The firm's assets are the project, $A=1.51$ \eqref{eq:1A}; $D=1$; and
\begin{equation}\label{eq:1E}
  E = A - D = 1.51 - 1 = 0.51 = C(F) = (3-1.785714)(0.42).
\end{equation}

\step{4}{Expected returns} Expected payoffs use the true probabilities $\tfrac12$:
\begin{align}
  r_A &= \frac{0.5(0.5)+0.5(3)}{1.51} - 1 = \frac{1.75}{1.51} - 1 = 0.158940, \label{eq:1rA}\\
  r_D &= \frac{0.5(0.5)+0.5(1.785714)}{1} - 1 = 1.142857 - 1 = 0.142857, \label{eq:1rD}\\
  r_E &= \frac{0.5(0) + 0.5(3-1.785714)}{0.51} - 1 = \frac{0.607143}{0.51} - 1 = 0.190476. \label{eq:1rE}
\end{align}

\step{5}{Check MM} The value-weighted average of \eqref{eq:1rD} and \eqref{eq:1rE} must equal
\eqref{eq:1rA}: $\tfrac{1}{1.51}(0.142857) + \tfrac{0.51}{1.51}(0.190476) = 0.094607 + 0.064333 = 0.158940$. \checkmark

\answer{$F=1.7857$; $A=1.51$, $D=1$, $E=0.51$; $r_A=15.89\%$, $r_D=14.29\%$, $r_E=19.05\%$.}

\remark{The promised yield, $F/D-1=78.57\%$, is far above the expected return on debt, 14.29\%:
creditors are promised 1.786 but receive only 0.5 in the slump. The yield is not the cost of debt.
The entrepreneur keeps the whole equity, worth $0.51$: exactly the project's NPV.}

\subsection*{(e$'$) Any mix: raise $D$ with debt and $1-D$ with equity}

\step{1}{Face value as a function of $D$} Debt is safe as long as $F\le0.5$, i.e.\ $D=0.92F\le0.46$.
\begin{equation}\label{eq:1Fgen}
  F(D) = \begin{cases}
    \dfrac{D}{0.92}, & 0 < D \le 0.46 \quad\text{(safe debt)},\\[8pt]
    \dfrac{D - (0.5)(0.50)}{0.42} = \dfrac{D-0.25}{0.42}, & 0.46 < D \le 1.51 \quad\text{(risky debt)}.
  \end{cases}
\end{equation}
The two branches meet at $D=0.46$, where both give $F=0.5$.

\step{2}{Values} $A=1.51$ for every $D$, and $E = 1.51 - D$. Outside shareholders pay $1-D$ for
their stake, so the entrepreneur's stake is worth $(1.51-D)-(1-D) = 0.51$, the NPV, for every mix.

\step{3}{Returns, safe region ($D\le0.46$)} $r_D=r_f=8.70\%$ and $r_A=15.89\%$ \eqref{eq:1rA}. MM~II
follows from $A\,r_A = D\,r_D + E\,r_E$:
\begin{equation}\label{eq:1mm2}
  r_E = r_A + \frac{D}{E}\,(r_A - r_D) .
\end{equation}
For example, $D=0.20$: $F = 0.20/0.92 = 0.217391$, $E=1.31$, and by \eqref{eq:1mm2}
\begin{equation}
  r_E = 0.158940 + \frac{0.20}{1.31}(0.071984) = 0.158940 + (0.152672)(0.071984) = 0.169930. \notag
\end{equation}

\step{4}{Returns, risky region ($D>0.46$)} Debt pays $(0.5,\ F)$, so
\begin{equation}\label{eq:1rDgen}
  r_D = \frac{0.5(0.5) + 0.5\,F(D)}{D} - 1 ,
\end{equation}
rising with $D$ (8.92\% at $D=0.47$, 14.29\% at $D=1$). Equity pays $(0,\ 3-F)$: nothing in the
slump. It is a multiple of the pure boom claim, so its expected return is that of the boom claim,
\begin{equation}\label{eq:1rEflat}
  r_E = \frac{0.5}{\psi_b} - 1 = \frac{0.5}{0.42} - 1 = 0.190476 ,
\end{equation}
the same for every $D$ in this region. Equation \eqref{eq:1mm2} still holds, now with the rising $r_D$.

\answer{$F(D)$ is given by \eqref{eq:1Fgen}; $A=1.51$ and $r_A=15.89\%$ for every $D$; $E=1.51-D$.
For $D\le0.46$: $r_D=8.70\%$ and $r_E$ rises linearly in $D/E$ from 15.89\% to 19.05\%. For
$D>0.46$: $r_D$ rises from 8.70\% to 14.29\% and $r_E$ stays at 19.05\%. The full table for
$D=0.01,\dots,1.00$ is in Appendix~A.}

\intuition{This is MM~II with risky debt. While debt is safe, each extra unit of leverage loads all
the asset risk onto equity, so $r_E$ rises linearly. Once debt defaults in the slump, creditors take
the slump payoff and start bearing risk. Equity is left with a boom-only claim whose risk can grow
no further, so $r_E$ flattens and the extra leverage shows up in a rising $r_D$. Through it all, the
WACC equals $r_A$ and the firm is worth 1.51.}

\subsection*{(f) Debt-financed dividend}

\step{1}{Face value} Raising 1.1 is in the risky region of \eqref{eq:1Fgen}:
$F = (1.1-0.25)/0.42 = 2.023810$.

\step{2}{Value of the firm} After the dividend, the firm holds only the project: $A = 1.51$, split as
$D = 1.1$ and $E = 1.51 - 1.1 = 0.41$. (Before the dividend it also held 0.1 of cash, so it was worth 1.61.)

\step{3}{The entrepreneur's wealth} Dividend plus equity: $0.1 + 0.41 = 0.51$, exactly what she had
with no dividend \eqref{eq:1E}.

\answer{The firm is worth 1.51 after the dividend ($D=1.1$, $E=0.41$). The entrepreneur is not better
off: her wealth stays at 0.51. The dividend only moves 0.1 of it from $t=1$ to $t=0$.}

\intuition{The creditors pay a fair 1.1 for a claim worth 1.1, so they transfer no value. The cash
paid out is the entrepreneur's own money, taken out early and paid for with a higher face value. She
could get the same $t=0$ consumption by selling 0.1 of her equity. With complete markets and no
frictions, payout and capital structure only change the timing of her wealth, not its amount. This
is the dividend-irrelevance argument of Topic 2B.}

% =============================================================================
\section{Exercise 2 --- Leverage, WACC and three debt policies}

\noindent\textbf{Data (millions of dollars).} 10 million shares, no debt; equity beta 1.20; FCF $=25$
per year in perpetuity from next year. Target $D/E = 0.25$, continuously rebalanced; $r_D=0.03$;
$\tau_c = 0.25$; CAPM with $r_f=0.03$ and $E[R_m]=0.08$.

\begin{notation}
  \item $r_U$, $\beta_U$: expected return and beta of the unlevered firm (its assets).
  \item $r_E$, $\beta_E$: of levered equity; $r_D=0.03$: cost of debt; $\tau_c=0.25$: tax rate.
  \item $V_U$, $V_L$: values of the unlevered and levered firm; $V_{TS}$: present value of the interest tax shields.
  \item $D$, $E$: market values of debt and equity; $N$: number of shares; $p$: share price.
\end{notation}

\subsection*{(a) Cost of equity before and after}

\step{1}{Before (all equity)} By the CAPM:
\begin{equation}\label{eq:2rU}
  r_U = r_f + \beta_U(E[R_m]-r_f) = 0.03 + 1.2(0.08-0.03) = 0.09 .
\end{equation}

\step{2}{After} With continuous rebalancing, debt moves with firm value, so the tax shields are as
risky as the assets: $r_{TS}=r_U$ (slides, Topic 2B). Then the pre-tax WACC equals $r_U$, and MM~II
keeps its no-tax form:
\begin{align}
  r_E &= r_U + \frac{D}{E}(r_U - r_D) \notag\\
      &= 0.09 + 0.25(0.09 - 0.03) = 0.105 . \label{eq:2rE}
\end{align}
Check through betas: $r_D=r_f$ means $\beta_D=0$, so $\beta_E = \beta_U(1+D/E) = 1.2\times1.25 = 1.5$
and $r_E = 0.03 + 1.5(0.05) = 0.105$.

\answer{$r_E = 9\%$ before and $10.5\%$ after.}

\subsection*{(b) New WACC}

$D/E=0.25$ means $D/V = 0.25/1.25 = 0.20$ and $E/V=0.80$. Using \eqref{eq:2rE}:
\begin{align}
  r_{WACC} &= \frac{E}{V}r_E + \frac{D}{V}r_D(1-\tau_c) \notag\\
           &= (0.80)(0.105) + (0.20)(0.03)(0.75) = 0.084 + 0.0045 = 0.0885 . \label{eq:2wacc}
\end{align}
Check with the slides' form for continuous rebalancing:
\begin{equation*}
  r_{WACC} = r_U - \frac{D}{V}\tau_c r_D = 0.09 - (0.20)(0.25)(0.03) = 0.0885 . \quad\checkmark
\end{equation*}

\answer{$r_{WACC} = 8.85\%$.}

\subsection*{(c) Values, price and repurchase}

\step{1}{Before} Using \eqref{eq:2rU}:
\begin{equation}\label{eq:2VU}
  V_U = \frac{25}{0.09} = 277.78, \qquad p_{\text{before}} = \frac{277.78}{10} = 27.78 .
\end{equation}

\step{2}{Levered firm and tax shield} Using \eqref{eq:2wacc}:
\begin{equation}\label{eq:2VL}
  V_L = \frac{25}{0.0885} = 282.49, \qquad V_{TS} = V_L - V_U = 282.49 - 277.78 = 4.71 .
\end{equation}

\step{3}{Price at the announcement} When the plan is announced, no debt has been issued yet. The
shareholders own the assets plus the expected tax shields, so the whole gain goes to them at once:
\begin{equation}\label{eq:2p}
  p = \frac{V_U + V_{TS}}{N} = \frac{282.49}{10} = 28.25 .
\end{equation}

\step{4}{Issue and repurchase} The debt issued is $D = 0.20\,V_L = 0.20(282.49) = 56.50$. It buys
back shares at the fair price \eqref{eq:2p}:
\begin{equation}\label{eq:2rep}
  \text{Shares repurchased} = \frac{56.50}{28.25} = 2.00, \qquad N_{\text{after}} = 10 - 2 = 8.00 .
\end{equation}
Check: $E = 0.80(282.49) = 225.99$ and $225.99/8 = 28.25$; the repurchase leaves the price unchanged.

\answer{$V_U = 277.78$; price 27.78 before, 28.25 after the announcement; $V_L = 282.49$, $V_{TS} = 4.71$;
2 million shares repurchased with \$56.50 million of debt; 8 million shares remain, at \$28.25.}

\subsection*{(d) Interest equal to 20\% of FCF}

Interest is $0.20\times25 = 5$ per year, so the tax shield is $\tau_c\times5 = 0.25\times5 = 1.25$ per
year. Because interest is tied to FCF, the shield is exactly as risky as FCF and is discounted at
$r_U$ \eqref{eq:2rU}:
\begin{equation}\label{eq:2TSd}
  V_{TS} = \frac{1.25}{0.09} = 13.89 .
\end{equation}

\answer{$V_{TS} = \$13.89$ million.}

\remark{Interest of 5 at $r_D=3\%$ corresponds to debt of about $5/0.03 = 166.67$ on average. That
is far more debt than in (c), hence the larger shield. The discount rate is still $r_U$, not $r_D$,
because the amount of debt moves with the firm's cash flows.}

\subsection*{(e) \$100 million of perpetual debt}

\step{1}{Tax shield} The debt is fixed and riskless, so the shield $\tau_c r_D D$ is discounted at $r_D$:
\begin{equation}\label{eq:2TSe}
  V_{TS} = \frac{\tau_c r_D D}{r_D} = \tau_c D = 0.25(100) = 25 .
\end{equation}

\step{2}{Price and shares} $V_L = 277.78 + 25 = 302.78$. At the announcement $p = 302.78/10 = 30.28$.
The repurchase buys $100/30.28 = 3.30$ million shares, leaving $10-3.30 = 6.70$ million. Check:
$E = 302.78 - 100 = 202.78$ and $202.78/6.70 = 30.28$.

\answer{$V_{TS} = \$25$ million; price \$30.28; 3.30 million shares repurchased; 6.70 million remain.}

\intuition{The three policies give shields of 4.71, 13.89 and 25. Two things drive the value: how
much debt there is, and how risky the shield is. A fixed amount of debt makes the shield as safe as
the debt; a policy tied to firm value or FCF makes it as risky as the assets. This is why the WACC
formula depends on the debt policy, and why APV is the safer tool when debt follows a known schedule.}

\refline{Slides Topic 2B (value of the tax shield; continuous rebalancing; WACC formulas and appendix);
Berk \S16.3 (the interest tax deduction and firm value; value of the interest tax shield; permanent debt; leverage and the WACC with taxes), pp.~473--478.}

% =============================================================================
\section{Exercise 3 --- Risky debt and who pays for bankruptcy costs}

\noindent\textbf{Data.} Investment $I=60$; FCF $=100$ or $60$ with probability $\tfrac12$ each;
$r_A = 0.15$; no taxes or other frictions.

\begin{notation}
  \item $V_U$, $V_L$: value of the all-equity and the levered firm; $D=60$: amount borrowed.
  \item $F$: face value of debt; $r_D$: expected return demanded by creditors; $E$: value of equity.
  \item $BC$: bankruptcy cost paid in default, a fraction $0.07$ of the realised cash flow.
\end{notation}

\subsection*{(a) All-equity firm}

\begin{equation}\label{eq:3VU}
  V_U = \frac{0.5(100) + 0.5(60)}{1.15} = \frac{80}{1.15} = 69.57 .
\end{equation}
The project's NPV is $69.57 - 60 = 9.57$.

\subsection*{(b) Debt of 60 at an expected return of 6.6\%}

\step{1}{Expected payoff creditors need} $60(1.066) = 63.96$.

\step{2}{Face value} If the debt were safe, $F = 63.96 > 60$, so it could not be paid in the bad
state: the debt is risky. Creditors get 60 (everything) in the bad state and $F$ in the good one:
\begin{equation}\label{eq:3Fb}
  0.5(60) + 0.5F = 63.96 \quad\Longrightarrow\quad F = 2(63.96) - 60 = 67.92 .
\end{equation}

\step{3}{Equity} By MM the levered firm is worth $V_U$, so
\begin{equation}\label{eq:3Eb}
  E = V_L - D = 69.57 - 60 = 9.57 .
\end{equation}
Equity pays $100 - 67.92 = 32.08$ in the good state and 0 in the bad state. Its expected payoff is
$16.04$, so $r_E = 16.04/9.57 - 1 = 67.69\%$. Check: $\tfrac{60}{69.57}(0.066) +
\tfrac{9.57}{69.57}(0.6769) = 0.15 = r_A$.

\answer{$F = 67.92$; $E = 9.57$ (the NPV, kept by the entrepreneur).}

\subsection*{(c) Bankruptcy costs of 7\% of the realised cash flow}

\step{1}{What creditors get in default} Default happens in the bad state ($60<F$). Lawyers take
$0.07(60) = 4.2$, so creditors receive $55.8$.

\step{2}{Face value} Keeping the creditors' required expected return at 6.6\%:
\begin{equation}\label{eq:3Fc}
  0.5(55.8) + 0.5F = 63.96 \quad\Longrightarrow\quad F = 2(63.96) - 55.8 = 72.12 .
\end{equation}

\step{3}{Value of the firm and of equity} Bankruptcy costs are a deadweight loss. The expected cost
is $0.5(4.2) = 2.1$. It is a part of the project's cash flow, so we discount it at $r_A$:
\begin{align}
  V_L &= V_U - PV(BC) = \frac{80 - 2.1}{1.15} = 69.57 - 1.83 = 67.74, \label{eq:3VLc}\\
  E   &= V_L - D = 67.74 - 60 = 7.74 . \label{eq:3Ec}
\end{align}

\step{4}{Who pays} Creditors still get a fair expected return on their 60; the higher face value
\eqref{eq:3Fc} compensates them for the 4.2 lost in default. Equity falls from 9.57 \eqref{eq:3Eb} to
7.74 \eqref{eq:3Ec}, a loss of $1.83 = PV(BC)$.

\answer{$F = 72.12$; $E = 7.74$. The shareholders (the entrepreneur) pay the whole present value of
the bankruptcy costs, 1.83. They pay it ex ante, through the higher face value, even though the
lawyers are paid out of the creditors' recovery in default.}

\intuition{Creditors anticipate the loss in default and price it in. The cost therefore falls on
whoever raises the money, at the time it is raised. This is why expected distress costs enter the
firm's capital-structure choice (the trade-off theory), even though shareholders get nothing in
default.}

\remark{A fully consistent pricing gives the same answer to the question ``who pays''. The market
data of (a) and (b) imply state prices $\psi_{good}=0.2982$ and $\psi_{bad}=0.6625$ (risk-free rate
4.10\%). The bad-state cost is then worth $4.2\times0.6625 = 2.78$, and equity is worth
$69.57 - 2.78 - 60 = 6.78$. Debt worth 60 would need $F = 77.25$, so the creditors' expected return
rises to 10.88\% because their claim is now riskier. The main answer keeps 6.6\% and discounts at
$r_A$, as the statement suggests. Either way, shareholders bear the full present value of the cost.}

\refline{Slides Topic 2B (bankruptcy costs; $V_L = V_U + V_{TS} - PV(\text{bankruptcy costs})$);
Berk \S16.4--16.5 (costs of bankruptcy and financial distress; the tradeoff theory), pp.~479--482.}

% =============================================================================
\section{Exercise 4 --- Yerba Industries: leverage, EPS and the P/E}

\noindent\textbf{Data.} All-equity, $\beta_U=1.2$, $r_U=0.125$. Issues risk-free debt at $0.05$ and
buys back 40\% of its shares. Perfect markets. Before: expected EPS \pounds1.50, forward P/E 14.

\begin{notation}
  \item $D/V=0.40$, $E/V=0.60$: weights after the transaction (the repurchase swaps 40\% of the
        equity for debt at market value, so $V$ does not change); $D/E = 0.4/0.6 = 2/3$.
  \item $\beta_D=0$, $r_D=0.05$: risk-free debt; $\beta_E$, $r_E$: of equity after the transaction.
\end{notation}

\subsection*{(a) Equity beta after}

With $\beta_D=0$, the asset beta is $\beta_U = \tfrac{E}{V}\beta_E$, so
\begin{equation}\label{eq:4beta}
  \beta_E = \beta_U\left(1+\frac{D}{E}\right) = 1.2\left(1+\frac{0.4}{0.6}\right) = 2.0 .
\end{equation}

\subsection*{(b) Expected return on equity after}

By MM~II:
\begin{equation}\label{eq:4rE}
  r_E = r_U + \frac{D}{E}(r_U - r_D) = 0.125 + \frac{2}{3}(0.125-0.05) = 0.175 .
\end{equation}
Check with the CAPM: the market risk premium is $(0.125-0.05)/1.2 = 0.0625$, so
$r_E = 0.05 + 2.0(0.0625) = 0.175$.

\subsection*{(c) Earnings per share after}

\step{1}{Price} Before: $p = 1.50\times14 = \pounds21$. Under MM the transaction does not change it.

\step{2}{Per original share} Debt raised per original share: $0.40\times21 = 8.40$. Interest:
$0.05\times8.40 = 0.42$. Earnings left for equity: $1.50 - 0.42 = 1.08$, now shared by
$0.60$ shares:
\begin{equation}\label{eq:4eps}
  EPS_{\text{after}} = \frac{1.08}{0.60} = \pounds1.80 .
\end{equation}

\answer{EPS rises from \pounds1.50 to \pounds1.80 (+20\%), but shareholders do not benefit: the share
price stays at \pounds21. The higher expected EPS is exactly the compensation for the higher risk of
leveraged equity ($\beta_E$ from 1.2 to 2.0, $r_E$ from 12.5\% to 17.5\%).}

\subsection*{(d) P/E after}

\begin{equation}\label{eq:4pe}
  P/E_{\text{after}} = \frac{21}{1.80} = 11.67 .
\end{equation}

\answer{The forward P/E falls from 14 to 11.67. This is reasonable: each pound of expected earnings
is now riskier, so investors pay less for it. The earnings yield rises from $1.50/21 = 7.14\%$ to
$1.80/21 = 8.57\%$, in line with the higher cost of equity.}

\intuition{Leverage raises EPS whenever the return on assets exceeds the cost of debt, and that says
nothing about value. Price is expected earnings capitalised at a rate that reflects their risk; MM~II
says the rate rises just enough to offset the higher EPS. A manager who targets EPS or ROE growth
through leverage creates no value in a perfect market (the slides' ``borrow more, value rises''
fallacy).}

\refline{Slides Topic 2A (MM II; leverage and EPS); Berk \S16.2 (the effect of leverage on risk and return; leverage and the cost of capital; capital structure fallacies), pp.~467--471.}

% =============================================================================
\appendix
\section{Exercise 1(e$'$): full table for $D = 0.01,\dots,1.00$}

Generated by \texttt{Listas/PS2-solutions\_check.py}, which also checks $\tfrac{D}{A}r_D +
\tfrac{E}{A}r_E = r_A$ on every row. Debt is safe up to $D=0.46$ ($F\le0.5$) and risky beyond.
Returns in \%.

{\footnotesize
\begin{longtable}{rrrrrrr}
  \toprule
  $D$ & $F$ & $r_D$ & $r_E$ & $r_A$ & $E$ & $A$ \\
  \midrule
  \endfirsthead
  \toprule
  $D$ & $F$ & $r_D$ & $r_E$ & $r_A$ & $E$ & $A$ \\
  \midrule
  \endhead
  \bottomrule
  \endfoot
  % Generated by PS2-solutions_check.py -- do not edit by hand.
0.01 & 0.0109 & 8.70 & 15.94 & 15.89 & 1.50 & 1.51 \\
0.02 & 0.0217 & 8.70 & 15.99 & 15.89 & 1.49 & 1.51 \\
0.03 & 0.0326 & 8.70 & 16.04 & 15.89 & 1.48 & 1.51 \\
0.04 & 0.0435 & 8.70 & 16.09 & 15.89 & 1.47 & 1.51 \\
0.05 & 0.0543 & 8.70 & 16.14 & 15.89 & 1.46 & 1.51 \\
0.06 & 0.0652 & 8.70 & 16.19 & 15.89 & 1.45 & 1.51 \\
0.07 & 0.0761 & 8.70 & 16.24 & 15.89 & 1.44 & 1.51 \\
0.08 & 0.0870 & 8.70 & 16.30 & 15.89 & 1.43 & 1.51 \\
0.09 & 0.0978 & 8.70 & 16.35 & 15.89 & 1.42 & 1.51 \\
0.10 & 0.1087 & 8.70 & 16.40 & 15.89 & 1.41 & 1.51 \\
0.11 & 0.1196 & 8.70 & 16.46 & 15.89 & 1.40 & 1.51 \\
0.12 & 0.1304 & 8.70 & 16.52 & 15.89 & 1.39 & 1.51 \\
0.13 & 0.1413 & 8.70 & 16.57 & 15.89 & 1.38 & 1.51 \\
0.14 & 0.1522 & 8.70 & 16.63 & 15.89 & 1.37 & 1.51 \\
0.15 & 0.1630 & 8.70 & 16.69 & 15.89 & 1.36 & 1.51 \\
0.16 & 0.1739 & 8.70 & 16.75 & 15.89 & 1.35 & 1.51 \\
0.17 & 0.1848 & 8.70 & 16.81 & 15.89 & 1.34 & 1.51 \\
0.18 & 0.1957 & 8.70 & 16.87 & 15.89 & 1.33 & 1.51 \\
0.19 & 0.2065 & 8.70 & 16.93 & 15.89 & 1.32 & 1.51 \\
0.20 & 0.2174 & 8.70 & 16.99 & 15.89 & 1.31 & 1.51 \\
0.21 & 0.2283 & 8.70 & 17.06 & 15.89 & 1.30 & 1.51 \\
0.22 & 0.2391 & 8.70 & 17.12 & 15.89 & 1.29 & 1.51 \\
0.23 & 0.2500 & 8.70 & 17.19 & 15.89 & 1.28 & 1.51 \\
0.24 & 0.2609 & 8.70 & 17.25 & 15.89 & 1.27 & 1.51 \\
0.25 & 0.2717 & 8.70 & 17.32 & 15.89 & 1.26 & 1.51 \\
0.26 & 0.2826 & 8.70 & 17.39 & 15.89 & 1.25 & 1.51 \\
0.27 & 0.2935 & 8.70 & 17.46 & 15.89 & 1.24 & 1.51 \\
0.28 & 0.3043 & 8.70 & 17.53 & 15.89 & 1.23 & 1.51 \\
0.29 & 0.3152 & 8.70 & 17.61 & 15.89 & 1.22 & 1.51 \\
0.30 & 0.3261 & 8.70 & 17.68 & 15.89 & 1.21 & 1.51 \\
0.31 & 0.3370 & 8.70 & 17.75 & 15.89 & 1.20 & 1.51 \\
0.32 & 0.3478 & 8.70 & 17.83 & 15.89 & 1.19 & 1.51 \\
0.33 & 0.3587 & 8.70 & 17.91 & 15.89 & 1.18 & 1.51 \\
0.34 & 0.3696 & 8.70 & 17.99 & 15.89 & 1.17 & 1.51 \\
0.35 & 0.3804 & 8.70 & 18.07 & 15.89 & 1.16 & 1.51 \\
0.36 & 0.3913 & 8.70 & 18.15 & 15.89 & 1.15 & 1.51 \\
0.37 & 0.4022 & 8.70 & 18.23 & 15.89 & 1.14 & 1.51 \\
0.38 & 0.4130 & 8.70 & 18.31 & 15.89 & 1.13 & 1.51 \\
0.39 & 0.4239 & 8.70 & 18.40 & 15.89 & 1.12 & 1.51 \\
0.40 & 0.4348 & 8.70 & 18.49 & 15.89 & 1.11 & 1.51 \\
0.41 & 0.4457 & 8.70 & 18.58 & 15.89 & 1.10 & 1.51 \\
0.42 & 0.4565 & 8.70 & 18.67 & 15.89 & 1.09 & 1.51 \\
0.43 & 0.4674 & 8.70 & 18.76 & 15.89 & 1.08 & 1.51 \\
0.44 & 0.4783 & 8.70 & 18.85 & 15.89 & 1.07 & 1.51 \\
0.45 & 0.4891 & 8.70 & 18.95 & 15.89 & 1.06 & 1.51 \\
0.46 & 0.5000 & 8.70 & 19.05 & 15.89 & 1.05 & 1.51 \\
0.47 & 0.5238 & 8.92 & 19.05 & 15.89 & 1.04 & 1.51 \\
0.48 & 0.5476 & 9.13 & 19.05 & 15.89 & 1.03 & 1.51 \\
0.49 & 0.5714 & 9.33 & 19.05 & 15.89 & 1.02 & 1.51 \\
0.50 & 0.5952 & 9.52 & 19.05 & 15.89 & 1.01 & 1.51 \\
0.51 & 0.6190 & 9.71 & 19.05 & 15.89 & 1.00 & 1.51 \\
0.52 & 0.6429 & 9.89 & 19.05 & 15.89 & 0.99 & 1.51 \\
0.53 & 0.6667 & 10.06 & 19.05 & 15.89 & 0.98 & 1.51 \\
0.54 & 0.6905 & 10.23 & 19.05 & 15.89 & 0.97 & 1.51 \\
0.55 & 0.7143 & 10.39 & 19.05 & 15.89 & 0.96 & 1.51 \\
0.56 & 0.7381 & 10.54 & 19.05 & 15.89 & 0.95 & 1.51 \\
0.57 & 0.7619 & 10.69 & 19.05 & 15.89 & 0.94 & 1.51 \\
0.58 & 0.7857 & 10.84 & 19.05 & 15.89 & 0.93 & 1.51 \\
0.59 & 0.8095 & 10.98 & 19.05 & 15.89 & 0.92 & 1.51 \\
0.60 & 0.8333 & 11.11 & 19.05 & 15.89 & 0.91 & 1.51 \\
0.61 & 0.8571 & 11.24 & 19.05 & 15.89 & 0.90 & 1.51 \\
0.62 & 0.8810 & 11.37 & 19.05 & 15.89 & 0.89 & 1.51 \\
0.63 & 0.9048 & 11.49 & 19.05 & 15.89 & 0.88 & 1.51 \\
0.64 & 0.9286 & 11.61 & 19.05 & 15.89 & 0.87 & 1.51 \\
0.65 & 0.9524 & 11.72 & 19.05 & 15.89 & 0.86 & 1.51 \\
0.66 & 0.9762 & 11.83 & 19.05 & 15.89 & 0.85 & 1.51 \\
0.67 & 1.0000 & 11.94 & 19.05 & 15.89 & 0.84 & 1.51 \\
0.68 & 1.0238 & 12.04 & 19.05 & 15.89 & 0.83 & 1.51 \\
0.69 & 1.0476 & 12.15 & 19.05 & 15.89 & 0.82 & 1.51 \\
0.70 & 1.0714 & 12.24 & 19.05 & 15.89 & 0.81 & 1.51 \\
0.71 & 1.0952 & 12.34 & 19.05 & 15.89 & 0.80 & 1.51 \\
0.72 & 1.1190 & 12.43 & 19.05 & 15.89 & 0.79 & 1.51 \\
0.73 & 1.1429 & 12.52 & 19.05 & 15.89 & 0.78 & 1.51 \\
0.74 & 1.1667 & 12.61 & 19.05 & 15.89 & 0.77 & 1.51 \\
0.75 & 1.1905 & 12.70 & 19.05 & 15.89 & 0.76 & 1.51 \\
0.76 & 1.2143 & 12.78 & 19.05 & 15.89 & 0.75 & 1.51 \\
0.77 & 1.2381 & 12.86 & 19.05 & 15.89 & 0.74 & 1.51 \\
0.78 & 1.2619 & 12.94 & 19.05 & 15.89 & 0.73 & 1.51 \\
0.79 & 1.2857 & 13.02 & 19.05 & 15.89 & 0.72 & 1.51 \\
0.80 & 1.3095 & 13.10 & 19.05 & 15.89 & 0.71 & 1.51 \\
0.81 & 1.3333 & 13.17 & 19.05 & 15.89 & 0.70 & 1.51 \\
0.82 & 1.3571 & 13.24 & 19.05 & 15.89 & 0.69 & 1.51 \\
0.83 & 1.3810 & 13.31 & 19.05 & 15.89 & 0.68 & 1.51 \\
0.84 & 1.4048 & 13.38 & 19.05 & 15.89 & 0.67 & 1.51 \\
0.85 & 1.4286 & 13.45 & 19.05 & 15.89 & 0.66 & 1.51 \\
0.86 & 1.4524 & 13.51 & 19.05 & 15.89 & 0.65 & 1.51 \\
0.87 & 1.4762 & 13.57 & 19.05 & 15.89 & 0.64 & 1.51 \\
0.88 & 1.5000 & 13.64 & 19.05 & 15.89 & 0.63 & 1.51 \\
0.89 & 1.5238 & 13.70 & 19.05 & 15.89 & 0.62 & 1.51 \\
0.90 & 1.5476 & 13.76 & 19.05 & 15.89 & 0.61 & 1.51 \\
0.91 & 1.5714 & 13.81 & 19.05 & 15.89 & 0.60 & 1.51 \\
0.92 & 1.5952 & 13.87 & 19.05 & 15.89 & 0.59 & 1.51 \\
0.93 & 1.6190 & 13.93 & 19.05 & 15.89 & 0.58 & 1.51 \\
0.94 & 1.6429 & 13.98 & 19.05 & 15.89 & 0.57 & 1.51 \\
0.95 & 1.6667 & 14.04 & 19.05 & 15.89 & 0.56 & 1.51 \\
0.96 & 1.6905 & 14.09 & 19.05 & 15.89 & 0.55 & 1.51 \\
0.97 & 1.7143 & 14.14 & 19.05 & 15.89 & 0.54 & 1.51 \\
0.98 & 1.7381 & 14.19 & 19.05 & 15.89 & 0.53 & 1.51 \\
0.99 & 1.7619 & 14.24 & 19.05 & 15.89 & 0.52 & 1.51 \\
1.00 & 1.7857 & 14.29 & 19.05 & 15.89 & 0.51 & 1.51 \\

\end{longtable}
}

\end{document}
